%% toc: P15 158 158.1,158.2,158.5,158.8
\section{Endpoint limits, Stieltjes laws and traces: proofs}
\label{sec:annex-p15-w02c-p15-003-stieltjes-endpoints}

Use the definitions of
Chapter~\ref{sec:p15-w02c-p15-003-stieltjes-endpoints},
\cref{def:p15-stieltjes-intervals} and \cref{def:p15-quantile-map}.
Throughout, \(F\) is a Stieltjes measure function. Since \(F\) is real
valued and nondecreasing, \(F(-\infty)\leq F(0)\leq F(\infty)\); hence
\(F(-\infty)\in[-\infty,\infty)\) and \(F(\infty)\in(-\infty,\infty]\).

\begin{annexitem}{annex:p15-w02c-p15-003-endpoint-limits}{thm:p15-w02c-p15-003-endpoint-limits}{Endpoint limits and truncation masses}
\begin{proof}
\paragraph{Part (1), the upper end.}
Suppose first that \(F(\infty)\) is finite, and let
\(\varepsilon>0\). Because \(F(\infty)\) is the least upper bound of the
values of \(F\), the smaller number \(F(\infty)-\varepsilon\) is not an
upper bound, so some real \(x_0\) has
\(F(x_0)>F(\infty)-\varepsilon\). For every \(x\geq x_0\),
monotonicity and the definition of the supremum give
\[
 F(\infty)-\varepsilon<F(x_0)\leq F(x)\leq F(\infty),
\]
so \(|F(x)-F(\infty)|<\varepsilon\). Suppose next that
\(F(\infty)=\infty\), and let \(M\) be real. The values of \(F\) are
not bounded above by \(M\), for otherwise their supremum would be at most
\(M\). So some real \(x_0\) has \(F(x_0)>M\), and for every
\(x\geq x_0\), \(F(x)\geq F(x_0)>M\). In both cases
\(F(x)\to F(\infty)\) as \(x\to\infty\).

\paragraph{Part (1), the lower end.}
Suppose first that \(F(-\infty)\) is finite, and let
\(\varepsilon>0\). Because \(F(-\infty)\) is the greatest lower bound
of the values of \(F\), the larger number
\(F(-\infty)+\varepsilon\) is not a lower bound, so some real \(x_0\)
has \(F(x_0)<F(-\infty)+\varepsilon\). For every \(x\leq x_0\),
\[
 F(-\infty)\leq F(x)\leq F(x_0)<F(-\infty)+\varepsilon ,
\]
so \(|F(x)-F(-\infty)|<\varepsilon\). Suppose next that
\(F(-\infty)=-\infty\), and let \(M\) be real. The values of \(F\) are
not bounded below by \(M\), so some real \(x_0\) has \(F(x_0)<M\), and
for every \(x\leq x_0\), \(F(x)\leq F(x_0)<M\). In both cases
\(F(x)\to F(-\infty)\) as \(x\to-\infty\).

\paragraph{Part (1), the integer sequences.}
Monotonicity gives \(F(n)\leq F(n+1)\) and \(F(-n-1)\leq F(-n)\).
Every \(F(n)\) is at most \(F(\infty)\), so
\(\sup_nF(n)\leq F(\infty)\). Conversely, every real \(x\) satisfies
\(x\leq n\) for some integer \(n\geq1\), by the Archimedean property, so
\(F(x)\leq F(n)\leq\sup_nF(n)\); taking the supremum over \(x\) gives
\(F(\infty)\leq\sup_nF(n)\). Likewise every \(F(-n)\) is at least
\(F(-\infty)\), and every real \(x\) satisfies \(-n\leq x\) for some
\(n\geq1\), so \(\inf_nF(-n)\leq F(-n)\leq F(x)\); taking the infimum
over \(x\) gives \(\inf_nF(-n)=F(-\infty)\). These hold whether or not
the endpoint values are finite.

\paragraph{Part (2), the shape of the truncations.}
A real \(x\) satisfies both \(a<x\) and \(-n<x\) exactly when
\(l_n=\max(a,-n)<x\), and satisfies both \(x\leq b\) and \(x\leq n\)
exactly when \(x\leq r_n=\min(b,n)\). Hence
\(K_n=\{x\in\mathbb R:l_n<x\leq r_n\}\), which is \((l_n,r_n]\) when
\(l_n<r_n\) and empty otherwise. As \(n\) increases, \(l_n\) does not
increase and \(r_n\) does not decrease, so once \(l_n<r_n\) holds it
holds for every larger \(n\). It holds for some \(n\): choose a real
\(c\) with \(a<c<b\), namely \((a+b)/2\) when both endpoints are finite,
\(b-1\) when only \(b\) is finite, \(a+1\) when only \(a\) is finite,
and \(0\) when neither is. For an integer \(n>|c|\), both \(a\) and
\(-n\) are below \(c\), and both \(b\) and \(n\) are above \(c\), so
\(l_n<c<r_n\). Let \(n_0\) be the least \(n\geq1\) with
\(l_n<r_n\). Then \(K_n=\varnothing\) and \(m_F(K_n)=0\) for
\(n<n_0\), while for \(n\geq n_0\)
\[
 m_F(K_n)=F(r_n)-F(l_n),
\]
a nonnegative real number. For \(n\geq n_0\),
\(F(r_n)\leq F(r_{n+1})\) and \(F(l_{n+1})\leq F(l_n)\), so
\(m_F(K_n)\leq m_F(K_{n+1})\); together with the zero values before
\(n_0\), the whole sequence of masses is nondecreasing.

\paragraph{Part (2), the endpoint values of the truncations.}
If \(b\) is finite, then \(r_n=b\) for every \(n\geq b\), so
\(F(r_n)=F(b)\) for all large \(n\). If \(b=\infty\), then
\(r_n=n\), and \(F(r_n)=F(n)\to F(\infty)\) by part~(1). If \(a\) is
finite, then \(l_n=a\) for every \(n\geq-a\), so \(F(l_n)=F(a)\) for
all large \(n\). If \(a=-\infty\), then \(l_n=-n\), and
\(F(l_n)=F(-n)\to F(-\infty)\) by part~(1). Thus in every case
\(F(r_n)\to F(b)\) and \(F(l_n)\to F(a)\), in the extended sense of
\cref{def:p15-w02c-p15-003-extended-limits} when the limit is infinite.

\paragraph{Part (2), the supremum identity.}
The value \(F(a)\) is never \(\infty\) and \(F(b)\) is never
\(-\infty\): for real endpoints they are real, and the remaining cases
are \(F(-\infty)\leq F(0)\) and \(F(\infty)\geq F(0)\). So if one of
the two values is infinite, then \(F(b)=\infty\) or
\(F(a)=-\infty\). Three cases therefore cover all possibilities.

\emph{Case A: \(F(a)\) and \(F(b)\) are finite.} The real sequences
\(F(r_n)\) and \(F(l_n)\) converge to \(F(b)\) and \(F(a)\), so their
differences \(m_F(K_n)\), for \(n\geq n_0\), converge to the real
number \(F(b)-F(a)\). A nondecreasing real sequence with a finite limit
has that limit as its supremum: no term exceeds the limit, because all
later terms would be at least that term, and for every
\(\varepsilon>0\) some term exceeds the limit minus \(\varepsilon\).
Hence \(\sup_nm_F(K_n)=F(b)-F(a)=m_F(I(a,b))\).

\emph{Case B: \(F(b)=\infty\).} Then \(b=\infty\), since \(F\) is
real at real points, and \(F(a)<\infty\). For \(m\geq n_0\),
\(l_m\leq l_{n_0}\), so
\[
 m_F(K_m)=F(m)-F(l_m)\geq F(m)-F(l_{n_0}).
\]
By part~(1), \(F(m)\) exceeds every real bound for large \(m\), and
\(F(l_{n_0})\) is a fixed real number, so \(m_F(K_m)\) exceeds every real
bound and \(\sup_nm_F(K_n)=\infty\). The assigned mass is
\(F(\infty)-F(a)\), which is \(\infty\) minus a finite number or minus
\(-\infty\), and therefore \(\infty\) by the convention of
\cref{def:p15-stieltjes-intervals}.

\emph{Case C: \(F(a)=-\infty\).} Then \(a=-\infty\) and
\(F(b)>-\infty\). For \(m\geq n_0\), \(r_m\geq r_{n_0}\), so
\[
 m_F(K_m)=F(r_m)-F(-m)\geq F(r_{n_0})-F(-m).
\]
By part~(1), \(F(-m)\) falls below every real bound for large \(m\),
and \(F(r_{n_0})\) is a fixed real number, so the masses exceed every
real bound and \(\sup_nm_F(K_n)=\infty\). The assigned mass is
\(F(b)-(-\infty)\), a finite number or \(\infty\) minus \(-\infty\),
which is \(\infty\) by the same convention.

When \(F(b)=\infty\) and \(F(a)=-\infty\) together, Cases~B and~C each
apply and give the same conclusion. In every case
\(\sup_nm_F(K_n)=m_F(I(a,b))\), and no infinite value has been
subtracted from another.
\end{proof}
\end{annexitem}

With this lemma, the truncation step in the proof of
\cref{thm:p15-stieltjes-construction} has each of its premises written:
the countable covering bound for bounded targets is applied to each
truncation \(K_n\), and the supremum identity of part~(2) carries it to
targets \(I(a,b)\) with \(a=-\infty\) or \(b=\infty\). That gives the
covering bound for all targets, the second partition condition, the
premeasure, and the interval masses \(\mu_F(I(a,\infty))=F(\infty)-F(a)\)
and \(\mu_F(\mathbb R)=F(\infty)-F(-\infty)\) recorded there.

\begin{annexitem}{annex:p15-w02c-p15-003-bounded-finite}{thm:p15-w02c-p15-003-bounded-finite}{Bounded sets and bounded intervals}
\begin{proof}
\paragraph{The bound.}
Let \(B\) be Borel with \(B\subseteq[-r,r]\), and let \(n>r\) be an
integer. Then \(-n<-r\), so \([-r,r]\subseteq(-n,n]=I(-n,n)\). By
\cref{thm:p15-stieltjes-construction},
\(\mu_F((-n,n])=F(n)-F(-n)\), a difference of real numbers.
Monotonicity of measures in \cref{thm:p15-measure-laws} gives
\(\mu_F(B)\leq F(n)-F(-n)<\infty\).

\paragraph{Bounded intervals are bounded Borel sets.}
A bounded interval is empty or has real endpoints \(c\leq d\). For real
\(c\),
\[
 \{c\}=\bigcap_{k\geq1}(c-1/k,c],
\]
since a real \(x\) lies in every \((c-1/k,c]\) exactly when
\(x\leq c\) and \(c-x<1/k\) for every \(k\), that is, when \(x=c\).
For \(c<d\),
\[
 (c,d)=\bigcup_{k\geq1}(c,d-1/k],
\]
where terms with \(d-1/k\leq c\) are empty: every \(x\) with
\(c<x<d\) satisfies \(x\leq d-1/k\) once \(1/k\leq d-x\), and no term
contains \(d\) or a point outside \((c,d)\). Further,
\([c,d]=\{c\}\cup(c,d]\) and \([c,d)=\{c\}\cup(c,d)\), each a disjoint
union. Every half-open interval \((c,d]\) is a member of
\(\mathcal S_{\mathbb R}\) and is Borel by
\cref{thm:p15-stieltjes-construction}; countable intersections and unions
of Borel sets are Borel. So every bounded interval is Borel, and it lies in
\([-r,r]\) for \(r=\max(|c|,|d|)\). The bound gives finiteness of its
measure.

\paragraph{The singleton.}
The sets \((c-1/k,c]\) decrease as \(k\) increases and have intersection
\(\{c\}\). The first has finite measure \(F(c)-F(c-1)\). Continuity from
above in \cref{thm:p15-measure-laws} gives
\[
 \mu_F(\{c\})=\lim_{k\to\infty}\bigl(F(c)-F(c-1/k)\bigr).
\]
The real sequence \(F(c-1/k)\) is nondecreasing in \(k\). Each term is at
most \(F(c-)\), because \(c-1/k<c\). For any \(t<c\), choose \(k\) with
\(1/k<c-t\); then \(t<c-1/k\) and \(F(t)\leq F(c-1/k)\). So the
supremum of the sequence is \(F(c-)\), and, being nondecreasing and
bounded above by \(F(c)\), the sequence converges to that supremum. Hence
\(\mu_F(\{c\})=F(c)-F(c-)\).

\paragraph{The other intervals.}
Let \(c<d\). The half-open case \(\mu_F((c,d])=F(d)-F(c)\) is the
defining interval mass. Since \((c,d]=(c,d)\cup\{d\}\) disjointly and all
measures involved are finite, \cref{thm:p15-measure-laws} gives
\[
 \begin{aligned}
 \mu_F((c,d))&=\mu_F((c,d])-\mu_F(\{d\})\\
   &=F(d)-F(c)-\bigl(F(d)-F(d-)\bigr)=F(d-)-F(c).
 \end{aligned}
\]
Finite additivity over the disjoint unions displayed above gives
\[
 \begin{aligned}
 \mu_F([c,d])&=\bigl(F(c)-F(c-)\bigr)+\bigl(F(d)-F(c)\bigr)=F(d)-F(c-),\\
 \mu_F([c,d))&=\bigl(F(c)-F(c-)\bigr)+\bigl(F(d-)-F(c)\bigr)=F(d-)-F(c-).
 \end{aligned}
\]
The degenerate bounded intervals are \([c,c]=\{c\}\), treated above, and
the empty interval, of measure zero.
\end{proof}
\end{annexitem}

\begin{annexitem}{annex:p15-w02c-p15-003-ray-laws}{thm:p15-w02c-p15-003-ray-laws}{Identity variables and closed rays}
\begin{proof}
\paragraph{Part (1).}
For every Borel set \(B\), \(\iota^{-1}(B)=B\in\mathcal B(\mathbb R)\).
So \(\iota\) is measurable from \((\mathbb R,\mathcal B(\mathbb R))\) to
\((\mathbb R,\mathcal B(\mathbb R))\), which on the probability space
\((\mathbb R,\mathcal B(\mathbb R),\nu)\) makes it a real random
variable in the sense of \cref{def:p15-measurable-map}. By
\cref{def:p15-pushforward-law}, its law is
\(B\mapsto\nu(\iota^{-1}(B))=\nu(B)\), which is \(\nu\). By
\cref{def:p15-cumulative-distribution}, its cumulative distribution
function is \(F_\iota(x)=\nu(\{\iota\leq x\})=\nu((-\infty,x])\).

\paragraph{Part (2), from the closed rays.}
Let \(\mathcal C=\{(-\infty,x]:x\in\mathbb R\}\). It is a pi-system,
since \((-\infty,x]\cap(-\infty,y]=(-\infty,\min(x,y)]\in\mathcal C\).
By \cref{thm:p15-borel-generators},
\(\sigma(\mathcal C)=\mathcal B(\mathbb R)\). The probability measures
\(\nu_1,\nu_2\) on \((\mathbb R,\sigma(\mathcal C))\) agree on every
member of \(\mathcal C\), so \cref{thm:p15-measure-uniqueness} gives
\(\nu_1(B)=\nu_2(B)\) for every Borel \(B\).

\paragraph{Part (2), from distribution functions.}
The same conclusion follows from \cref{thm:p15-cdf-laws}. By part~(1),
\(\nu_i\) is the law of the identity variable \(\iota_i\) on
\((\mathbb R,\mathcal B(\mathbb R),\nu_i)\), whose cumulative
distribution function is \(x\mapsto\nu_i((-\infty,x])\). These two
functions are equal by hypothesis. The last assertion of
\cref{thm:p15-cdf-laws} says that the two real random variables
\(\iota_1,\iota_2\) then have the same law, that is, \(\nu_1=\nu_2\).

\paragraph{Part (3), the measure \(\mu_F\).}
The function \(F\) is real valued, nondecreasing and right continuous,
so it is a Stieltjes measure function. By part~(1) of
\cref{thm:p15-w02c-p15-003-endpoint-limits}, \(F(x)\to F(-\infty)\) as
\(x\to-\infty\); by hypothesis \(F(x)\to0\). If \(F(-\infty)\) were
finite and nonzero, then with \(\varepsilon=|F(-\infty)|/2\) all
sufficiently negative \(x\) would satisfy both
\(|F(x)-F(-\infty)|<\varepsilon\) and \(|F(x)|<\varepsilon\), which
is impossible because the distance from \(F(-\infty)\) to zero is
\(2\varepsilon\). If \(F(-\infty)=-\infty\), all sufficiently negative
\(x\) would satisfy both \(F(x)<-1\) and \(|F(x)|<1\), again
impossible. So \(F(-\infty)=0\). At the upper end, \(F(x)\to F(\infty)\)
and \(F(x)\to1\) as \(x\to\infty\). If \(F(\infty)\) were finite and
different from one, the same choice
\(\varepsilon=|F(\infty)-1|/2\) gives a contradiction; if
\(F(\infty)=\infty\), all sufficiently large \(x\) would satisfy both
\(F(x)>2\) and \(|F(x)-1|<1\), which is impossible. So
\(F(\infty)=1\). By
\cref{thm:p15-stieltjes-construction},
\[
 \begin{aligned}
 \mu_F(\mathbb R)&=F(\infty)-F(-\infty)=1,\\
 \mu_F((-\infty,x])&=\mu_F(I(-\infty,x))=F(x)-F(-\infty)=F(x).
 \end{aligned}
\]
So \(\mu_F\) is a probability measure on
\((\mathbb R,\mathcal B(\mathbb R))\). By part~(1), it is the law of
the identity variable on \((\mathbb R,\mathcal B(\mathbb R),\mu_F)\),
whose cumulative distribution function is
\(x\mapsto\mu_F((-\infty,x])=F(x)\).

\paragraph{Part (3), the law of \(Q_F\).}
\Cref{thm:p15-cdf-existence} states that \(Q_F\) is a real random
variable on \((\Omega_U,\mathcal F_U,\mathbb U)\) and that
\(Q_F(u)\leq x\) exactly when \(u\leq F(x)\). In
\cref{annex:p15.cdf-existence} these two conclusions are proved from the
definition of \(Q_F\), right continuity and the ray generator criterion,
before and without the comparison of laws; only they are used here. For
real \(x\), the ray equivalence gives
\[
 \{u\in(0,1):Q_F(u)\leq x\}=\{u\in(0,1):u\leq F(x)\}=\{U\leq F(x)\},
\]
where \(U(u)=u\) is the identity variable of
\cref{thm:p15-lebesgue-unit-interval}. With \(t=F(x)\in[0,1]\), that
theorem gives \(\mathbb U(U\leq t)\) equal to \(0\) when \(t=0\), to
\(t\) when \(0<t<1\), and to \(1\) when \(t=1\); in every case it equals
\(F(x)\). By \cref{thm:p15-pushforward-laws}, the law \(\mathbb U_Q\)
of \(Q_F\) is a probability measure on \((\mathbb R,\mathcal B(\mathbb R))\),
and
\[
 \mathbb U_Q((-\infty,x])=\mathbb U(Q_F\leq x)=F(x)=\mu_F((-\infty,x])
\]
for every real \(x\). Part~(2) gives \(\mathbb U_Q=\mu_F\).

\paragraph{Part (3), uniqueness.}
If a probability measure \(\nu\) on \((\mathbb R,\mathcal B(\mathbb R))\)
satisfies \(\nu((-\infty,x])=F(x)\) for every real \(x\), then \(\nu\)
and \(\mu_F\) agree on every closed ray, and part~(2) gives
\(\nu=\mu_F\). Conversely \(\mu_F\) has these ray values. Let \(X\) be a
real random variable on a probability space \((\Omega,\mathcal F,P)\).
Its law \(P_X\) is a probability measure on
\((\mathbb R,\mathcal B(\mathbb R))\) by \cref{thm:p15-pushforward-laws},
and \(P_X((-\infty,x])=P(X\leq x)=F_X(x)\). If \(F_X=F\), the preceding
sentence applied to \(\nu=P_X\) gives \(P_X=\mu_F\). If \(P_X=\mu_F\),
then \(F_X(x)=\mu_F((-\infty,x])=F(x)\) for every \(x\).
\end{proof}
\end{annexitem}

The quantile and uniqueness conclusions of \cref{thm:p15-cdf-existence}
are therefore part~(3) of \cref{thm:p15-w02c-p15-003-ray-laws}: the law of
\(Q_F\) is \(\mu_F\), and \(\mu_F\) is the only Borel probability law,
whether presented as a measure with given ray values or as the law of a
real random variable, whose cumulative distribution function is \(F\).

\begin{annexitem}{annex:p15-w02c-p15-003-trace-generation}{thm:p15-w02c-p15-003-trace-generation}{Traces and generators}
\begin{proof}
Use \cref{def:p15-w02c-p15-003-trace-sigma-field}. Since \(D\) is
nonempty, \cref{def:p15-sigma-algebra} and \cref{thm:p15-generated-events}
apply with \(D\) in place of \(\Omega\); in particular \(\sigma_D(\mathcal K)\)
is a sigma-algebra on \(D\) containing \(\mathcal K\) and contained in every
sigma-algebra on \(D\) that contains \(\mathcal K\), for every
\(\mathcal K\subseteq2^D\).

\paragraph{Part (1).}
Every member \(A\cap D\) of \(\mathcal F_D\) is a subset of \(D\).
The whole space \(D\) is \(\Omega\cap D\), and \(\Omega\in\mathcal F\),
so \(D\in\mathcal F_D\). Let \(B=A\cap D\) with \(A\in\mathcal F\). Its
complement relative to \(D\) is
\[
 D\setminus B=\{\omega\in D:\omega\notin A\}=(\Omega\setminus A)\cap D,
\]
and \(\Omega\setminus A\in\mathcal F\), so \(D\setminus B\in\mathcal F_D\).
Let \(B_k=A_k\cap D\) with \(A_k\in\mathcal F\) for each \(k\) in a
countable index set. Then
\[
 \bigcup_kB_k=\Bigl(\bigcup_kA_k\Bigr)\cap D,
\]
because a point lies in the left side exactly when it lies in \(D\) and in
some \(A_k\). The union of the \(A_k\) lies in \(\mathcal F\), so the
union of the \(B_k\) lies in \(\mathcal F_D\). Hence \(\mathcal F_D\) is
a sigma-algebra on \(D\).

\paragraph{Part (2).}
Let \(D\in\mathcal F\). For \(A\in\mathcal F\), the set \(A\cap D\) is a
countable intersection of two events, so it lies in \(\mathcal F\) by
\cref{thm:p15-generated-events}, and it is a subset of \(D\). Conversely,
if \(A\in\mathcal F\) and \(A\subseteq D\), then \(A=A\cap D\in\mathcal F_D\).
This proves \(\mathcal F_D=\{A\in\mathcal F:A\subseteq D\}\); in
particular \(\mathcal F_D\subseteq\mathcal F\), so \(\mu|_D\) is defined
on all of \(\mathcal F_D\) and takes values in \([0,\infty]\).
Its value at \(\varnothing=\varnothing\cap D\) is \(\mu(\varnothing)=0\).
If \((B_n)_{n\geq1}\) are pairwise disjoint members of \(\mathcal F_D\),
they are pairwise disjoint members of \(\mathcal F\), their union lies in
\(\mathcal F_D\) by part~(1), and countable additivity of \(\mu\) gives
\[
 \mu|_D\Bigl(\bigcup_nB_n\Bigr)=\mu\Bigl(\bigcup_nB_n\Bigr)
       =\sum_n\mu(B_n)=\sum_n\mu|_D(B_n).
\]
So \(\mu|_D\) is a measure on \((D,\mathcal F_D)\) by
\cref{def:p15-measure}. Its total mass is \(\mu|_D(D)=\mu(D)\), so it is
finite when \(\mu(D)<\infty\).

\paragraph{Part (3), first inclusion.}
Let \(\mathcal F=\sigma(\mathcal C)\), put
\(\mathcal C_D\defeq\{C\cap D:C\in\mathcal C\}\subseteq2^D\), and
\(\mathcal G\defeq\sigma_D(\mathcal C_D)\). Every \(C\in\mathcal C\)
lies in \(\mathcal F\), so \(C\cap D\in\mathcal F_D\); thus
\(\mathcal C_D\subseteq\mathcal F_D\). By part~(1), \(\mathcal F_D\) is a
sigma-algebra on \(D\) containing \(\mathcal C_D\). The minimality of
\(\sigma_D(\mathcal C_D)\) among such sigma-algebras gives
\(\mathcal G\subseteq\mathcal F_D\).

\paragraph{Part (3), second inclusion.}
Let
\[
 \mathcal H\defeq\{A\subseteq\Omega:A\cap D\in\mathcal G\}.
\]
It contains \(\mathcal C\), since \(C\cap D\in\mathcal C_D\subseteq\mathcal G\)
for \(C\in\mathcal C\). It is a sigma-algebra on \(\Omega\). First,
\(\Omega\cap D=D\), and \(D\in\mathcal G\) because \(\mathcal G\) is a
sigma-algebra on \(D\); so \(\Omega\in\mathcal H\). Second, if
\(A\in\mathcal H\), then
\[
 (\Omega\setminus A)\cap D=D\setminus(A\cap D),
\]
which is the complement of \(A\cap D\in\mathcal G\) relative to \(D\) and so
lies in \(\mathcal G\); hence \(\Omega\setminus A\in\mathcal H\). Third, if
\(A_k\in\mathcal H\) for each \(k\) in a countable index set, then
\((\bigcup_kA_k)\cap D=\bigcup_k(A_k\cap D)\), as in part~(1), is a
countable union of members of \(\mathcal G\) and lies in \(\mathcal G\);
hence \(\bigcup_kA_k\in\mathcal H\). By \cref{thm:p15-generated-events},
\(\mathcal F=\sigma(\mathcal C)\subseteq\mathcal H\). That is, for every
\(A\in\mathcal F\), the set \(A\cap D\) lies in \(\mathcal G\), which says
\(\mathcal F_D\subseteq\mathcal G\).

The two inclusions give \(\mathcal F_D=\sigma_D(\mathcal C_D)\). For
\(\mathcal C=\{E_j:j\geq1\}\), the collection \(\mathcal C_D\) is
\(\{E_j\cap D:j\geq1\}\), and the identity reads
\(\mathcal F_D=\sigma_D(E_1\cap D,E_2\cap D,\ldots)\). Neither inclusion
used \(D\in\mathcal F\).

\paragraph{Part (4).}
Let \(B\in\mathcal S\). Each \(\omega\in\Omega\) lies in exactly one
\(D_i\), which is then nonempty, and \(g(\omega)=g_i(\omega)\). Hence
\[
 g^{-1}(B)=\bigcup_{i\in I,\ D_i\neq\varnothing}g_i^{-1}(B),
\]
where \(g_i^{-1}(B)=\{\omega\in D_i:g_i(\omega)\in B\}\). Measurability
of \(g_i\) from \((D_i,\mathcal F_{D_i})\) gives
\(g_i^{-1}(B)\in\mathcal F_{D_i}\), and part~(2), applied to
\(D_i\in\mathcal F\), gives \(\mathcal F_{D_i}\subseteq\mathcal F\).
The union is countable because \(I\) is, so \(g^{-1}(B)\in\mathcal F\)
for every \(B\in\mathcal S\). By \cref{def:p15-measurable-map}, \(g\) is
measurable.
\end{proof}
\end{annexitem}

\begin{annexitem}{annex:p15-w02c-p15-003-block-data}{thm:p15-w02c-p15-003-block-data}{Blocks of finite positive mass}
\begin{proof}
Since \(\nu(\varnothing)=0<\nu(D)\), the set \(D\) is nonempty, so
\cref{thm:p15-w02c-p15-003-trace-generation} applies to it.

Part~(1) is parts~(1) and~(3) of
\cref{thm:p15-w02c-p15-003-trace-generation} with
\(\mathcal C=\{E_j:j\geq1\}\).

For part~(2), part~(2) of \cref{thm:p15-w02c-p15-003-trace-generation}
makes \(\mu|_D\) and \(\nu|_D\) measures on \((D,\mathcal F_D)\), with
\(\mu|_D\) finite because \(\mu(D)<\infty\). Every value \(\nu|_D(B)\),
\(B\in\mathcal F_D\), satisfies
\(\nu|_D(B)\leq\nu|_D(B)+\nu|_D(D\setminus B)=\nu(D)<\infty\), by
additivity over the disjoint sets \(B\) and \(D\setminus B\) together with
infinitely many empty sets. So \(\lambda(B)=\nu(B)/\nu(D)\) is a real number
in \([0,1]\), and \(\lambda(\varnothing)=0\), \(\lambda(D)=1\). For
pairwise disjoint \(B_n\in\mathcal F_D\), the partial sums of
\(\nu|_D(B_n)\) are real, are bounded by \(\nu(D)\), and increase to
\(\nu|_D(\bigcup_nB_n)\); dividing each partial sum by the positive real
\(\nu(D)\) preserves this, so the partial sums of \(\lambda(B_n)\)
increase to \(\lambda(\bigcup_nB_n)\). By \cref{def:p15-measure},
\(\lambda\) is a probability measure on \((D,\mathcal F_D)\).

For part~(3), let \(B\in\mathcal F_D\) with \(\lambda(B)=0\). Then
\(\nu(B)=\nu(D)\lambda(B)=0\). By part~(2) of
\cref{thm:p15-w02c-p15-003-trace-generation}, \(B\in\mathcal F\), so
\(\mu\ll\nu\) gives \(\mu(B)=0\), that is \(\mu|_D(B)=0\). This is
\(\mu|_D\ll\lambda\) in the sense of \cref{def:p15-absolute-continuity}
on \((D,\mathcal F_D)\).

For part~(4), \(A\cap D\in\mathcal F_D\) by the definition of the trace,
and \(\mu|_D(A\cap D)=\mu(A\cap D)\) by the definition of the
restriction.

Parts~(1)--(3) are the hypotheses of the finite case of
\cref{thm:p15-countably-generated-radon-nikodym}, read on the space \(D\):
a sigma-field generated by the sequence \((E_j\cap D)_{j\geq1}\), a finite
measure \(\mu|_D\), a probability measure \(\lambda\), and
\(\mu|_D\ll\lambda\).
\end{proof}
\end{annexitem}

\begin{annexitem}{annex:p15-w02c-p15-003-block-integrals}{thm:p15-w02c-p15-003-block-integrals}{Integrals on a block and on the whole space}
\begin{proof}
Throughout, \(D\in\mathcal F\) is nonempty, so part~(2) of
\cref{thm:p15-w02c-p15-003-trace-generation} gives
\(\mathcal F_D=\{A\in\mathcal F:A\subseteq D\}\subseteq\mathcal F\) and
makes \(\nu|_D\) a measure on \((D,\mathcal F_D)\). Measurability of a
function \(h\) with values in \([0,\infty]\subseteq\overline{\mathbb R}\)
is tested by \cref{thm:p15-measurability-criterion} on the generators
\([-\infty,a]\), \(a\in\mathbb R\), of
\cref{def:p15-extended-random-variables}: \(h\) is measurable exactly when
every set \(\{h\leq a\}\) belongs to the sigma-algebra of its domain.

\paragraph{Part (1), measurability.}
Let \(a\in\mathbb R\). If \(a<0\), then \(\{\bar g\leq a\}=\varnothing\),
since \(\bar g\geq0\). If \(a\geq0\), then \(\bar g\leq a\) holds at every
point of \(\Omega\setminus D\), where \(\bar g=0\), and at a point of \(D\)
exactly when \(g\leq a\) there; so
\[
 \{\bar g\leq a\}=\{\omega\in D:g(\omega)\leq a\}\cup(\Omega\setminus D).
\]
The first set lies in \(\mathcal F_D\subseteq\mathcal F\) by measurability
of \(g\), and the second lies in \(\mathcal F\) because \(D\) does. Hence
\(\bar g\) is measurable from \((\Omega,\mathcal F)\).

\paragraph{Part (1), simple functions.}
Let \(s\) be a nonnegative simple function on \((D,\mathcal F_D)\), written
as in \cref{def:p15-simple-integral} as
\(s=\sum_{j=1}^ma_j\mathbf1_{B_j}\) with finite \(a_j\geq0\) and pairwise
disjoint \(B_j\in\mathcal F_D\) whose union is \(D\). The sets
\(B_1,\ldots,B_m,\Omega\setminus D\) are pairwise disjoint members of
\(\mathcal F\) with union \(\Omega\), and
\(\bar s=\sum_{j=1}^ma_j\mathbf1_{B_j}+0\cdot\mathbf1_{\Omega\setminus D}\)
is a representation of the zero extension of \(s\) as a nonnegative simple
function on \((\Omega,\mathcal F)\). With the convention
\(0\cdot\infty=0\) of \cref{def:p15-simple-integral},
\[
 \int_\Omega\bar s\,d\nu
  =\sum_{j=1}^ma_j\nu(B_j)+0\cdot\nu(\Omega\setminus D)
  =\sum_{j=1}^ma_j\,\nu|_D(B_j)
  =\int_Ds\,d(\nu|_D).
\]
The simple integrals do not depend on the representation, by
\cref{thm:p15-simple-integral-laws}.

Conversely, let \(t\) be a nonnegative simple function on
\((\Omega,\mathcal F)\) with \(t\leq\bar g\), written as
\(t=\sum_{j=1}^ma_j\mathbf1_{A_j}\) with pairwise disjoint
\(A_j\in\mathcal F\) whose union is \(\Omega\). On \(\Omega\setminus D\),
\(0\leq t\leq\bar g=0\), so \(t=0\) there. The sets \(A_j\cap D\) are
pairwise disjoint members of \(\mathcal F_D\) with union \(D\), so the
restriction \(t_D=\sum_{j=1}^ma_j\mathbf1_{A_j\cap D}\) of \(t\) to \(D\) is
a nonnegative simple function on \((D,\mathcal F_D)\), and \(t_D\leq g\)
because \(t\leq\bar g\) and \(\bar g=g\) on \(D\). Since \(t\) vanishes
off \(D\), it is the zero extension of \(t_D\), and the display gives
\(\int_\Omega t\,d\nu=\int_Dt_D\,d(\nu|_D)\).

\paragraph{Part (1), the integrals.}
Let \(L_D\) be the set of numbers \(\int_Ds\,d(\nu|_D)\) for nonnegative
simple \(s\) on \((D,\mathcal F_D)\) with \(s\leq g\), and let
\(L_\Omega\) be the set of numbers \(\int_\Omega t\,d\nu\) for nonnegative
simple \(t\) on \((\Omega,\mathcal F)\) with \(t\leq\bar g\). If \(s\leq g\)
on \(D\), then \(\bar s\leq\bar g\) on \(\Omega\), since both equal \(0\)
off \(D\); so \(\int_Ds\,d(\nu|_D)=\int_\Omega\bar s\,d\nu\) lies in
\(L_\Omega\). If \(t\leq\bar g\), then
\(\int_\Omega t\,d\nu=\int_Dt_D\,d(\nu|_D)\) lies in \(L_D\). Hence
\(L_D=L_\Omega\), and their suprema, which are \(\int_Dg\,d(\nu|_D)\) and
\(\int_\Omega\bar g\,d\nu\) by \cref{def:p15-nonnegative-integral}, are
equal.

\paragraph{Part (2), scaling.}
Let \(0<\nu(D)<\infty\) and \(c=1/\nu(D)\), a positive real number. For
\(B\in\mathcal F_D\), additivity of \(\nu|_D\) over the disjoint sets \(B\)
and \(D\setminus B\) gives \(\nu(B)\leq\nu(D)<\infty\), so
\(\lambda(B)=c\,\nu(B)\) is a real number in \([0,1]\). Moreover
\(\lambda(\varnothing)=0\), and for pairwise disjoint
\(B_n\in\mathcal F_D\) the real partial sums of \(\nu(B_n)\) increase to
\(\nu(\bigcup_nB_n)\); multiplying them by \(c\) shows that the partial sums
of \(\lambda(B_n)\) increase to \(\lambda(\bigcup_nB_n)\). So \(\lambda\) is
a measure, indeed a probability measure, on \((D,\mathcal F_D)\).

Let \(h:D\to[0,\infty]\) be measurable from \((D,\mathcal F_D)\). For a
nonnegative simple \(s=\sum_ja_j\mathbf1_{B_j}\) on \((D,\mathcal F_D)\),
all the numbers \(a_j\nu(B_j)\) are real, and
\[
 \int_Ds\,d\lambda=\sum_ja_j\,c\,\nu(B_j)=c\sum_ja_j\nu(B_j)
  =c\int_Ds\,d(\nu|_D).
\]
Whether \(s\leq h\) does not depend on the measure, so the set of simple
integrals that defines \(\int_Dh\,d\lambda\) is the image under
\(x\mapsto cx\) of the set that defines \(\int_Dh\,d(\nu|_D)\). This map is
an increasing bijection of \([0,\infty]\) onto itself, with
\(c\cdot\infty=\infty\), so it carries the supremum of a set to the
supremum of its image. Therefore
\(\int_Dh\,d\lambda=c\int_Dh\,d(\nu|_D)\), and by part~(1),
\[
 \int_Dh\,d\lambda=c\int_\Omega\bar h\,d\nu.
\]

\paragraph{Part (2), integrals over a member of the trace.}
Let \(B\in\mathcal F_D\); then \(B\in\mathcal F\). Let \(h\) equal \(g\) on
\(B\) and \(0\) on \(D\setminus B\), so that \(\int_Bg\,d\lambda\) is
\(\int_Dh\,d\lambda\). The function \(h\) is measurable from
\((D,\mathcal F_D)\): for \(a<0\) the set \(\{h\leq a\}\) is empty, and for
\(a\geq0\) it is \((\{\omega\in D:g(\omega)\leq a\}\cap B)\cup(D\setminus B)\),
a member of \(\mathcal F_D\) by part~(1) of
\cref{thm:p15-w02c-p15-003-trace-generation}. The zero extension
\(\bar h\) equals \(g\) on \(B\) and \(0\) elsewhere in \(\Omega\), which is
the function \(\bar g\mathbf1_B\). The scaling identity gives
\[
 \int_Bg\,d\lambda=\int_Dh\,d\lambda=c\int_\Omega\bar g\mathbf1_B\,d\nu
  =\frac1{\nu(D)}\int_B\bar g\,d\nu .
\]

\paragraph{Part (3).}
For \(N\in\mathcal F_D\), \(\lambda(N)=c\,\nu(N)\) with \(c>0\), so
\(\lambda(N)=0\) exactly when \(\nu(N)=0\). Now let \(g<\infty\) hold
\(\lambda\)-almost everywhere: by \cref{def:p15-nonnegative-integral}
there is \(M\in\mathcal F_D\) with \(\lambda(M)=0\) and \(g<\infty\) on
\(D\setminus M\). The set
\[
 N_\infty\defeq\{\omega\in D:g(\omega)=\infty\}
  =D\setminus\bigcup_{n\geq1}\{\omega\in D:g(\omega)\leq n\}
\]
lies in \(\mathcal F_D\) by measurability of \(g\) and part~(1) of
\cref{thm:p15-w02c-p15-003-trace-generation}, and \(N_\infty\subseteq M\).
Monotonicity, part~(1) of \cref{thm:p15-measure-laws}, gives
\(\lambda(N_\infty)\leq\lambda(M)=0\), hence \(\nu(N_\infty)=0\); and
\(N_\infty\in\mathcal F\) because \(\mathcal F_D\subseteq\mathcal F\).
\end{proof}
\end{annexitem}

\begin{annexitem}{annex:p15-w02c-p15-003-assembly}{thm:p15-w02c-p15-003-assembly}{Assembly of local densities}
\begin{proof}
For \(i\in I^+\) put \(c_i=1/\nu(D_i)\), a positive real number, so that
\(f=c_if_i\) on \(D_i\). Every \(i\in I^+\) has \(D_i\neq\varnothing\),
because \(\nu(\varnothing)=0\).

\paragraph{Measurability.}
Measurability is tested on the sets \(\{h\leq a\}\), \(a\in\mathbb R\), as
in the proof of \cref{thm:p15-w02c-p15-003-block-integrals}. For
\(i\in I^+\) and real \(a\), the sets \(\{\omega\in D_i:c_if_i(\omega)\leq a\}\)
and \(\{\omega\in D_i:f_i(\omega)\leq a\,\nu(D_i)\}\) coincide: at a point
where \(f_i\) is real, multiplying by the positive real \(c_i\) or by its
reciprocal preserves the inequality, and at a point where \(f_i=\infty\)
both \(c_if_i=\infty>a\) and \(f_i>a\,\nu(D_i)\). The second set lies in
\(\mathcal F_{D_i}\), so \(c_if_i\) is measurable from
\((D_i,\mathcal F_{D_i})\). For \(i\notin I^+\) with \(D_i\neq\varnothing\),
the zero function on \(D_i\) is measurable, since each set
\(\{0\leq a\}\) is \(\varnothing\) or \(D_i\). Part~(4) of
\cref{thm:p15-w02c-p15-003-trace-generation}, with
\(S=\overline{\mathbb R}\) and the sigma-algebra of
\cref{def:p15-extended-random-variables}, shows that \(f\) is measurable
from \((\Omega,\mathcal F)\); its values lie in \([0,\infty]\).

\paragraph{The identity on one block.}
Fix \(A\in\mathcal F\) and \(i\in I\). The set \(A\cap D_i\) lies in
\(\mathcal F\) by \cref{thm:p15-generated-events}.

Let \(i\in I^+\). Then \(B\defeq A\cap D_i\in\mathcal F_{D_i}\), and by
hypothesis \(\int_Bf_i\,d\lambda_i=\mu(B)\). Part~(2) of
\cref{thm:p15-w02c-p15-003-block-integrals}, with \(D=D_i\),
\(\lambda=\lambda_i\) and \(g=f_i\), gives
\[
 \mu(B)=\int_Bf_i\,d\lambda_i=c_i\int_\Omega\bar f_i\mathbf1_B\,d\nu
       =\int_\Omega c_i\bar f_i\mathbf1_B\,d\nu,
\]
the last step by the scalar rule of
\cref{thm:p15-nonnegative-integral-laws} with the finite constant
\(c_i\geq0\). The functions \(c_i\bar f_i\mathbf1_B\) and
\(f\mathbf1_B\) are equal at every point: on \(B\subseteq D_i\) both are
\(c_if_i\), with \(c_i\cdot\infty=\infty\) where \(f_i=\infty\), and off
\(B\) both are \(0\). Hence
\(\int_{A\cap D_i}f\,d\nu=\int_\Omega f\mathbf1_B\,d\nu=\mu(A\cap D_i)\).

Let \(i\notin I^+\), so \(\nu(D_i)=0\). Monotonicity, part~(1) of
\cref{thm:p15-measure-laws}, gives \(\nu(A\cap D_i)=0\), and \(\mu\ll\nu\)
gives \(\mu(A\cap D_i)=0\). Since \(f=0\) on \(D_i\), the function
\(f\mathbf1_{A\cap D_i}\) is the zero function, whose integral is
\(0\cdot\nu(\Omega)=0\) by \cref{def:p15-simple-integral} and
\cref{thm:p15-monotone-integral}. So
\(\int_{A\cap D_i}f\,d\nu=0=\mu(A\cap D_i)\) in this case as well.

\paragraph{Summation over the blocks.}
By part~(4) of \cref{thm:p15-nonnegative-integral-laws}, the function
\(E\mapsto\int_Ef\,d\nu\) is a measure on \((\Omega,\mathcal F)\). The sets
\(A\cap D_i\), \(i\in I\), are pairwise disjoint members of \(\mathcal F\)
with union \(A\), and the index set is countable; a finite family is
completed by empty sets, which contribute \(0\) to both sides. Countable
additivity of this measure and of \(\mu\) gives
\[
 \int_Af\,d\nu=\sum_{i\in I}\int_{A\cap D_i}f\,d\nu
              =\sum_{i\in I}\mu(A\cap D_i)=\mu(A).
\]

\paragraph{Finiteness almost everywhere.}
For \(i\in I^+\) let \(N_i=\{\omega\in D_i:f_i(\omega)=\infty\}\). By
part~(3) of \cref{thm:p15-w02c-p15-003-block-integrals}, applied with
\(D=D_i\) and \(g=f_i\), which is finite \(\lambda_i\)-almost everywhere,
each \(N_i\) is a member of \(\mathcal F\) with \(\nu(N_i)=0\). On \(D_i\)
with \(i\in I^+\), \(f=c_if_i\) is infinite exactly where \(f_i\) is,
because \(c_i\) is a positive real; on \(D_i\) with \(i\notin I^+\), \(f=0\).
Hence \(\{f=\infty\}=\bigcup_{i\in I^+}N_i\), a countable union of
\(\nu\)-null members of \(\mathcal F\), which is \(\nu\)-null by part~(2) of
\cref{thm:p15-measure-laws}. So \(f<\infty\) holds \(\nu\)-almost everywhere.

\paragraph{Application to the countably generated case.}
In the sigma-finite case of \cref{thm:p15-countably-generated-radon-nikodym},
the blocks \(D_i\) are members of \(\mathcal F\) with
\(\mu(D_i)<\infty\) and \(\nu(D_i)<\infty\), pairwise disjoint with union
\(\Omega\). For \(i\in I^+\), \cref{thm:p15-w02c-p15-003-block-data}
shows that the finite case applies on \((D_i,\mathcal F_{D_i})\) to the
finite measure \(\mu|_{D_i}\), the probability \(\lambda_i\) and the
generators \(E_j\cap D_i\). Its conclusion is a nonnegative function
\(f_i\), measurable for \(\mathcal F_{D_i}\) and finite
\(\lambda_i\)-almost everywhere, with
\(\int_Bf_i\,d\lambda_i=\mu|_{D_i}(B)=\mu(B)\) for every
\(B\in\mathcal F_{D_i}\). These are the hypotheses of the present theorem,
whose conclusions are the representation
\(\int_Af\,d\nu=\mu(A)\) for every \(A\in\mathcal F\) and the finiteness of
\(f\) outside a \(\nu\)-null set.
\end{proof}
\end{annexitem}
